% =====================================================================
%  Supreme: swapping the memory of a world model
%  Build:  pdflatex supreme.tex   (twice, for references)
%  No external .bib -- references are inlined.
% =====================================================================
\documentclass[11pt,a4paper]{article}

\usepackage[margin=1in]{geometry}
\usepackage{amsmath,amssymb}
\usepackage{booktabs}
\usepackage{graphicx}
\usepackage{xcolor}
\usepackage[colorlinks=true,linkcolor=blue!50!black,urlcolor=blue!50!black,
            citecolor=blue!50!black]{hyperref}
\usepackage{caption}
\usepackage{enumitem}
\usepackage{listings}

\lstset{basicstyle=\ttfamily\footnotesize,frame=single,breaklines=true,
        columns=flexible,showstringspaces=false,xleftmargin=1em}

\newcommand{\code}[1]{\texttt{\small #1}}
\newcommand{\WM}{\textsc{World Model}}
\newcommand{\SUP}{\textsc{Supreme}}
\newcommand{\todo}[1]{\textcolor{red}{[#1]}}

\title{\SUP: Isolating the Memory Layer of a World Model\\
       \large An apples-to-apples comparison of the MDN-RNN and a
       nested-learning memory, at laptop scale}
\author{ Kuppala Sarath Narendra , Neelakandan NC}
\date{\today}

\begin{document}
\maketitle

\begin{abstract}
We build a world model in the sense of Ha and Schmidhuber
\cite{ha2018world}---a visual encoder $V$, a memory $M$, and a controller
$C$---and replace \emph{only} the recurrent core of $M$ with the memory
mechanism proposed by Behrouz et al.\ under the nested-learning framework
\cite{behrouz2025nested}, specifically the \textsc{Hope} module: a
self-modifying Titans block followed by a Continuum Memory System. Every other
component---the environments, the data pipeline, the convolutional VAE, the
mixture-density output head, the loss, the optimiser, the controller, and the
entire evaluation suite---is literally the same code in both arms, so the
comparison has exactly one axis of variation. Because the only thing that
differs is the memory, we evaluate with memory benchmarks: cue retention as a
function of delay, prediction accuracy at the one transition per episode that
requires memory, open-loop imagination divergence, and end-to-end control.

At a 32-step required memory horizon \textsc{Hope} solves the task outright
(decision accuracy $1.000$, real-environment return $1.00$) while the MDN-RNN and
a GRU control both sit at chance, and a linear probe locates the difference
precisely: \textsc{Hope}'s state still encodes the cue at delay 32, the GRU's has
decayed to chance, the LSTM's is gone by delay 8. Two ablations sharpen the
claim rather than support it. Attaching the Continuum Memory System to the
baseline LSTM changes nothing at all, so the advantage lies entirely in the
sequence mixer; and replacing the self-modifying block with softmax attention
also solves the task, so self-modification is not uniquely necessary for
retention. What distinguishes the two is imagination: attention retains
perfectly yet its open-loop rollout diverges ($44.6$ MSE at horizon 16 against
$0.51$), while \textsc{Hope} both retains and stays on trajectory at a
fixed-size state. We also quantify a property of the objective itself---
maximum-likelihood world-model training rewards retaining a task-critical cue by
only ${\sim}0.7$ nats out of ${\sim}25$, so at the standard weighting every
memory layer we tested scores at chance. The entire study---data collection,
training, evolutionary search and evaluation---runs on a single Apple M4 laptop
with 16\,GB of unified memory.
\end{abstract}

\tableofcontents

% =====================================================================
\section{Motivation and scope}

The claim under test is narrow and should stay narrow:

\begin{quote}
Swapping the recurrent core of a world model's memory module for a
nested-learning memory changes how much the world model remembers.
\end{quote}

\noindent The answer turns out to be yes, but for a narrower reason than the
framing suggests, and \S\ref{sec:ablation} is where that narrowing happens.

Two consequences follow immediately, and they shape everything below.

\paragraph{One axis of variation, enforced by code layout.} The obvious
implementation---a self-contained copy of the pipeline per model---cannot
support this claim. With two copies of the data loader, the VAE, the loss and
the controller, the two arms drift, and every reported difference becomes ``the
memory layer, plus whatever else diverged''; no diff of two sixty-file trees
tells you which. We therefore keep a single shared core and reduce each model to
one small folder containing its recurrent core and nothing else
(\S\ref{sec:layout}). A model folder \emph{cannot} change the training loop
because it contains no training loop.

\paragraph{Memory benchmarks, not control benchmarks.} If the memory is the only
difference, then a benchmark that does not stress memory cannot resolve the
question. We show in \S\ref{sec:findings-nll} that the standard world-model
metric---one-step negative log-likelihood---separates an LSTM core from a GRU
core by $0.01$ nats while their effective memory horizons differ threefold. Our
primary instruments are therefore probes of the recurrent state and prediction
accuracy restricted to memory-critical transitions.

% =====================================================================
\section{The two models}
\label{sec:models}

Both models have the three-part structure of \cite{ha2018world}:
\begin{equation}
\text{obs}_t \;\xrightarrow{\;V\;}\; z_t
\qquad
a_t = C\!\left([z_t\,;h_t]\right)
\qquad
h_{t+1} = M(z_t, a_t, h_t).
\end{equation}
$V$ compresses a frame to a latent, $M$ predicts the next latent, and $C$ maps
the pair to an action.

\subsection{Shared components}

These are identical---the same objects, not merely the same hyperparameters---in
both arms.

\paragraph{$V$: convolutional VAE.} A faithful reimplementation of
\cite[\S2.1]{ha2018world}: four strided convolutions with channel widths
$(32,64,128,256)$ to a $1024$-d bottleneck, a diagonal Gaussian latent of $32$
dimensions, and a mirrored deconvolutional decoder producing $64\times64\times3$
frames. This is $4{,}348{,}547$ parameters, matching the count reported in their
appendix. The objective is squared reconstruction error plus a KL term with a
per-dimension floor (free bits, $0.5$ nats). The floor matters here: the
synthetic memory environments contain very low-entropy frames, and an
unconstrained KL collapses most latent dimensions to the prior, which would
render the cue unrecoverable from $z$ and cripple \emph{both} memory models.

$V$ is trained once per environment and then \emph{frozen}; latents are
precomputed once and both arms consume the same \code{latents.npz}. Retraining
$V$ per model would let VAE seed noise leak into the memory comparison, and that
noise is larger than the effect being measured. It also makes the study
affordable: after $V$ is frozen, training $M$ touches a ${\sim}30$\,MB array of
32-d vectors rather than a multi-gigabyte frame store.

\paragraph{Output head: mixture density network.} $M$ is a backbone plus a
\emph{shared} head. The head is a diagonal MDN with $K=5$ components per latent
dimension, exactly as in \cite[\S2.2]{ha2018world}, with optional scalar reward
and termination heads. Writing $\pi,\mu,\sigma$ for the head's outputs, the
per-step loss is
\begin{equation}
\mathcal{L}_t \;=\; -\sum_{d=1}^{32}\log \sum_{k=1}^{5}
   \pi_{t,d,k}\,\mathcal{N}\!\left(z_{t+1,d};\,\mu_{t,d,k},\,\sigma_{t,d,k}\right).
\label{eq:mdn}
\end{equation}
Sharing the head is not a convenience. If one arm had an MDN head and the other
a Gaussian head, a difference in likelihood would confound \emph{what the model
remembers} with \emph{how it expresses uncertainty}, and no ablation could
separate them. Fixing the head makes a likelihood difference attributable to the
recurrent core by construction.

\paragraph{$C$: linear controller.} A single affine map from $[z_t;h_t]$ to an
action, trained by CMA-ES \cite{hansen2006cma} (with a dependency-free
Gaussian-ES fallback). Ha and Schmidhuber keep $C$ trivial so that competence
must live in $V$ and $M$; here the argument does double duty. Because $C$ is
linear and identical in both arms, a difference in return is a difference in how
much task-relevant information the memory state exposes \emph{linearly}---which
is precisely what our probes measure. Setting \code{controller.hidden\_input =
false} gives the no-memory ablation in which $C$ sees $z_t$ only.

\subsection{Model A --- \WM{} (baseline)}

The control condition is the recurrent half of the MDN-RNN
\cite[\S2.2]{ha2018world}: a single-layer LSTM with $256$ hidden units consuming
$[z_t;a_t]$, whose output feeds the shared MDN head. Nothing else. With a 32-d
latent and a 2-d one-hot action this is $299{,}008$ backbone parameters and
$422{,}882$ for the full $M$.

\subsection{Model B --- \SUP{} (\textsc{Hope})}

The challenger is the \textsc{Hope} module of \cite[\S8.3]{behrouz2025nested}: a
self-modifying Titans block whose output passes through a Continuum Memory
System chain.

\paragraph{Self-modifying Titans (\cite[\S8.1--8.2]{behrouz2025nested}, Eqs.\
86--93).} Rather than projecting the input to keys, values and queries with
weights frozen after pre-training, every projection is itself an associative
memory that keeps learning during the sequence, and---the self-modifying
part---each memory generates its own regression targets:
\begin{align}
q_t &= x_t W_q, &
k_t &= \mathcal{M}_{k,t-1}(x_t), &
v_t &= \mathcal{M}_{v,t-1}(x_t), \\
\eta_t &= \mathcal{M}_{\eta,t-1}(x_t), &
\alpha_t &= \mathcal{M}_{\alpha,t-1}(x_t), &
\hat v_{\square,t} &= \mathcal{M}_{\square,t-1}(v_t), \\
\mathcal{M}_{\square,t} &= \mathcal{M}_{\square,t-1}
   \left(\alpha_t I - \eta_t k_t k_t^{\!\top}\right)
   - \eta_t \nabla \mathcal{L}\!\left(\mathcal{M}_{\square,t-1};k_t,\hat v_{\square,t}\right),
   & o_t &= \mathcal{M}_{\mathrm{mem},t-1}(q_t)
\label{eq:titans}
\end{align}
for $\square\in\{k,v,\eta,\alpha,\mathrm{mem}\}$, with $\mathcal{L}$ the $L_2$
regression loss, $\eta_t$ a data-dependent inner learning rate and $\alpha_t$ a
data-dependent retention gate. The initial states $\mathcal{M}_{\square,0}$ are
ordinary trainable parameters, meta-learned across sequences. We use the
matrix-valued memories for which \cite[Eq.~93]{behrouz2025nested} derives the
closed form, retaining the residual retrieval structure of their Eq.~89
($\mathcal{M}(x) = x + Mx$), so the inner gradient is exact without invoking
autograd inside the recurrence. Following \cite[\S8.3]{behrouz2025nested} we
$L_2$-normalise $q$ and $k$ and apply a depthwise causal convolution of window~4.

The chunked schedule of \cite[\S8.2]{behrouz2025nested} is implemented
faithfully rather than as an optimisation: keys, values, gates, generated
targets and inner gradients are all computed from the memory state at the
\emph{start} of the chunk, which changes what the block computes.
\code{chunk\_size = 1} recovers the fully online form.

\paragraph{Continuum Memory System (\cite[\S7.1]{behrouz2025nested}, Eqs.\
70--71).} A chain of MLP blocks in which level $\ell$ is updated only every
$C^{(\ell)}$ steps, from the gradient accumulated over that window:
\begin{equation}
y_t = \mathrm{MLP}^{(f_k)}\!\left(\cdots \mathrm{MLP}^{(f_1)}(o_t)\right),
\qquad
\theta^{(f_\ell)}_{i+1} =
\begin{cases}
\theta^{(f_\ell)}_i - \sum_{t=i-C^{(\ell)}}^{i}\eta\,\nabla\mathcal{L}(\theta_t;x_t)
   & i \equiv 0 \ (\mathrm{mod}\ C^{(\ell)}),\\[2pt]
\theta^{(f_\ell)}_i & \text{otherwise.}
\end{cases}
\label{eq:cms}
\end{equation}
The argument is that when a fast level is overwritten by new data, the knowledge
it held survives in the slower levels, so the system forgets reluctantly. We use
periods $(1,4,16)$ and the sequential variant of their Eq.~73.

\paragraph{Composition.} $\textsc{Hope} = \text{Titans} \rightarrow \text{CMS}$,
their Eqs.~94--97. Because \textsc{Hope} \emph{is} the composition of two
separable mechanisms, we implement them as independently switchable parts, which
yields the ablation table for free (\S\ref{sec:variants}).

\subsection{Controls and ablations}
\label{sec:variants}

\begin{table}[h]
\centering
\begin{tabular}{llr}
\toprule
Key & Composition & Backbone params \\
\midrule
\code{lstm}           & LSTM(256) --- the MDN-RNN core \cite{ha2018world} & $299{,}008$ \\
\code{gru}            & GRU(256) --- gated-RNN control & $224{,}256$ \\
\code{attention}      & causal self-attention --- ``perfect memory'' control & $299{,}058$ \\
\midrule
\code{supreme}        & Titans $\rightarrow$ CMS (\textsc{Hope}) & $300{,}696$ \\
\code{supreme-titans} & Titans only & $224{,}520$ \\
\code{supreme-cms}    & LSTM $\rightarrow$ CMS & $375{,}184$ \\
\code{hope-attention} & attention $\rightarrow$ CMS (their \textsc{Hope-Attention}) & $375{,}234$ \\
\bottomrule
\end{tabular}
\caption{The seven memory cores. Only \code{supreme} is parameter-matched to the
baseline, and only it needs to be---it is the model making the claim.
\code{supreme-titans} and \code{supreme-cms} are a strict subset and a strict
superset of it by construction, so their counts are reported rather than
equalised.}
\label{tab:backbones}
\end{table}

The attention control deserves comment. Behrouz et al.\ characterise softmax
attention as a ``perfect memory'': a non-parametric solution that caches every
past token and never compresses \cite[\S8]{behrouz2025nested}. It is therefore
the strongest available retention baseline, and a memory paper comparing only
against recurrent cores leaves the obvious question unanswered. The interesting
comparison is not who retains more---attention retains everything by
construction---but at what price: attention's state grows linearly with the
episode and its compute quadratically, while a fast-weight core keeps a
fixed-size state.

\subsection{What differs, and what cannot}
\label{sec:layout}

\begin{lstlisting}
World Model/          MODEL 1: lstm_backbone.py, attention_backbone.py
Supreme/              MODEL 2: titans.py, cms.py, hope.py
wmcore/               SHARED: envs, data, V, MDN head, loss, C, dream, benchmarks
\end{lstlisting}

Model folders are discovered by path rather than by import name, which is why one
can be called \code{World Model} despite the space. A machine check,
\code{assert\_comparable}, diffs two run configurations and aborts unless the only
differing keys are the backbone identifier and its keyword arguments; it runs
before training in every study.

\subsection{Deviations from the source papers}
\label{sec:deviations}

Recorded in full because each is a place where a reproduction could legitimately
differ.

\begin{enumerate}[leftmargin=*]
\item \textbf{Is $q$ produced by a memory?} \cite[Eq.~94]{behrouz2025nested}
lists a memory for $q$, while the surrounding text states that $q_t=x_tW_q$ is
``the only non-adaptive projection''. We follow the text.

\item \textbf{What is $\mathcal{L}$ in Eq.~71?} The equation states it is ``the
objective of choice for the task at hand'', i.e.\ the \emph{outer} task loss.
Read literally, Eq.~71 therefore describes a gradient-accumulation schedule over
optimiser steps, and that is what we implement. This is consistent with their M3
optimiser (\S7.2), which applies the same idea to momentum across optimiser
steps, and with \S7.3's proposal to initialise CMS blocks from pre-trained MLP
weights. The paper also gestures at a within-sequence reading; the two are not
equivalent and we took the unambiguous one.

\item \textbf{Causality in Eq.~90.} The retrieval memory is indexed
$C\lceil t/C\rceil$, which is not causal for $t$ inside a chunk. We use the
chunk start.

\item \textbf{Memory-norm projection.} Written exactly as Eq.~88 the recurrence
diverges---measured here, the memory norm reaches $10^{10}$ by step 80 and
overflows by step 96. The mechanism is a property of self-modification rather
than an implementation error: expanding the update, the generated target
$\hat v=\mathcal{M}_{\mathrm{ref}}(v)$ is itself a function of the memory and
$v$ is in turn $\mathcal{M}_{\mathrm{ref}}(x)$, so the write term contains
$+\eta\,\mathcal{M}_{\mathrm{ref}}\,v\,k^{\!\top}$, which is \emph{quadratic} in
$\lVert M\rVert$. In the paper the memories are two-layer MLPs inside a
normalised stack that absorbs this; the linear form has nothing to absorb it. We
project each head's memory onto a Frobenius ball after each update, leaving the
update rule exactly as Eq.~88. At the study's operating point the projection
never engages ($\lVert M\rVert_F\!\approx\!2.5$ against a radius of $20$); it
engages only on rollouts long enough for the divergence to appear, and the
engagement rate is logged every epoch.

Two less invasive-looking alternatives were tried first and are recorded because
both are tempting and both are worse. Normalising the generated target $\hat v$
stabilises the recurrence but destroys the magnitude of the quantity the update
regresses on: measured, decision accuracy at corridor~32 fell from $1.000$ to
chance while the training curve looked healthy for the first ${\sim}25$ epochs,
so the damage was invisible early. Tying retention to the write rate
($\alpha\le1-\eta$) is harmless but pointless---the learned $\alpha$ settles near
$0.77$ either way, so the constraint never binds.

\item \textbf{Scale.} Everything here is three orders of magnitude smaller than
\cite{behrouz2025nested}. Their claims are about language models; ours are about
a $300$\,k-parameter memory inside a world model. Nothing in this report should
be read as a reproduction of their results.
\end{enumerate}

% =====================================================================
\section{Benchmarks}
\label{sec:bench}

\subsection{Environments}

\begin{table}[h]
\centering
\begin{tabular}{lll}
\toprule
Environment & Tests & Role \\
\midrule
\code{TMaze-v0}          & retention of one cue across a tunable delay & primary \\
\code{SequenceRecall-v0} & holding $k$ associations at once            & capacity axis \\
\code{CarRacing-v0}      & the reference task of \cite{ha2018world}    & external comparability \\
\bottomrule
\end{tabular}
\caption{Environments. All emit $64\times64\times3$ uint8 frames.}
\end{table}

The two synthetic environments are pure NumPy---no gymnasium, Box2D or
pygame---so they install instantly and are bit-reproducible. Three properties
make them suitable for a world-model study specifically.

\emph{The memory requirement lands inside the world-model loss, not only in the
reward.} In \code{TMaze-v0} a coloured cue appears at $t=0$, a corridor of
tunable length follows with visually salient but uninformative distractors, the
agent chooses at a junction, and the next frame reveals whether it was right.
Predicting that one frame requires both the cue (seen at $t=0$) and the action
(taken now); nothing in the current frame reveals it. The environment tags that
transition \code{memory\_critical} and the benchmark reports likelihood
restricted to exactly those steps.

\emph{Ground-truth latent variables are emitted in \code{info}.} The cue, the
correct action and the delay since the cue are exposed to the evaluation code and
never to the model, which makes retention curves exact rather than inferred.

\emph{The environment is verified as a memory task rather than assumed to be
one.} A unit test confirms an agent given the cue scores $40/40$ while a
cue-blind agent is at chance.

Corridor length is the study's independent variable: it is the required memory
horizon, and it can be swept in minutes.

\subsection{Metrics}

\begin{description}[leftmargin=*,style=nextline]
\item[Decision accuracy (headline).] At the memory-critical transition, we take
the model's predicted next latent, assign it to the nearer of the two
ground-truth outcome centroids, and compare with the truth. Chance is exactly
$0.5$, and---crucially---it cannot be gamed by hedging: a model that predicts the
average of the two outcomes lands between the centroids and scores at chance.
Hedging is precisely the likelihood-optimal strategy for a model that has not
learned the rule, which is why likelihood alone is the wrong instrument.

\item[Teacher-forced NLL.] Equation~\ref{eq:mdn} over all steps, and separately
restricted to memory-critical steps.

\item[Open-loop imagination.] Warm the state on ground truth, then roll forward
feeding the model its own predictions with the true action sequence. Error at
small $k$ measures local dynamics; the slope measures whether the state retains
enough to stay on the true trajectory.

\item[Retention probes.] Harvest the recurrent features $h_t$, and for each
cue-to-probe delay $d$ fit a \emph{linear} ridge classifier from $h_t$ to the
ground-truth cue, splitting \emph{by episode}. Reported alongside is the control:
the same probe on $z_t$ alone. Linear probes rather than MLPs, because a
nonlinear probe can manufacture accuracy from a state no linear controller could
use, and $C$ here \emph{is} linear.

\item[Continual learning.] Task-incremental protocol over environment variants
that share pixel statistics and differ only in the cue$\rightarrow$outcome rule,
with no replay and no task label, reported as backward and forward transfer in
the style of Lopez-Paz and Ranzato \cite{lopezpaz2017gem}. Forgetting is only
measurable against something learned, so the benchmark self-reports
\code{informative: false} when mean accuracy on each task immediately after
training it is near chance.

\item[Control.] Real-environment return with the CMA-ES controller, plus the
dream-transfer variant of \cite[\S5]{ha2018world} in which $C$ is trained
entirely inside $M$'s imagination and only evaluated for real.

\item[Cost.] Parameters, per-step inference latency, training-pass time,
throughput and peak resident memory, on every run.
\end{description}

\subsection{Validating the instruments}

Seven of the twenty-five tests in the repository exist to validate the
measurement apparatus rather than the code:

\begin{itemize}[leftmargin=*]
\item the open-loop harness is checked against an \emph{analytic oracle}---a
model that is exactly right must score zero error at every horizon, and a biased
model must compound. This test caught a defect in which the first open-loop step
was silently teacher-forced;
\item decision accuracy is checked with both an oracle (must score $1.0$) and a
hedger (must score chance);
\item the probe is checked to find signal in a planted-signal dataset and to
\emph{reject} it in pure noise, and to lose accuracy when the group split is
honest, which pins the episode-level split;
\item the MDN likelihood is checked against the closed-form Gaussian;
\item every backbone must satisfy $\text{forward(sequence)} =
\text{step-by-step}$ to $10^{-5}$. A core failing this trains normally,
benchmarks normally on teacher-forced metrics, and is silently wrong in every
open-loop, dream and control result.
\end{itemize}

\subsection{A property of the objective that the benchmark must correct for}
\label{sec:findings-nll}

A delayed-recall episode contains exactly one transition that cannot be predicted
without memory---once the answer is revealed, memory is no longer needed. At
sequence length $L$ that transition carries ${\sim}1/L$ of the loss, and the most
a model can gain by getting it right is the entropy of the outcome: about $0.69$
nats for a binary cue, against a per-step NLL of ${\sim}25$.

The consequence is measurable. On \code{TMaze-v0} with corridor 6, sweeping the
weight $\lambda$ applied to memory-critical transitions:

\begin{table}[h]
\centering
\begin{tabular}{lcccc}
\toprule
$\lambda$ & Decision acc.\ (chance $0.5$) & NLL @ critical & NLL (all steps) & Probe at decision \\
\midrule
1   & $0.500$ & $25.53$ & $13.16$ & $0.44$ (chance) \\
20  & $1.000$ & $19.42$ & $12.76$ & $1.00$ \\
100 & $1.000$ & $13.68$ & $12.49$ & $1.00$ \\
\bottomrule
\end{tabular}
\caption{LSTM core, 2000 rollouts, 60 epochs. At $\lambda=1$ the converged model
predicts almost exactly the \emph{midpoint} of the two possible successors---mean
distance to the midpoint $0.25$ against a centroid separation of $6.74$---while a
probe shows the cue was held in the state four steps into the corridor and gone
by the decision.}
\label{tab:lambda}
\end{table}

The model is not failing to remember. It is correctly optimising an objective
that barely rewards remembering, and spending its recurrent capacity on the many
locally predictable corridor frames instead. Two things follow. First, this is
worth reporting on its own: maximum-likelihood world-model training gives a
recurrent core almost no incentive to retain task-critical information, and the
incentive can be quantified. Second, the benchmark has no headroom at
$\lambda=1$---every memory layer scores at chance and nothing can be
distinguished---so the memory configurations use $\lambda=20$, applied
identically to both arms and disclosed. Note that it also \emph{lowers}
unweighted NLL, so it is a fix to an optimisation-priority pathology rather than
a quality trade-off.

% =====================================================================
\section{Machine and software environment}
\label{sec:env}

Every result in this report was produced on a single machine.

\begin{table}[h]
\centering
\begin{tabular}{ll}
\toprule
\multicolumn{2}{l}{\textbf{Hardware}}\\
\midrule
Chip                & Apple M4 \\
CPU                 & 10 cores (4 performance, 6 efficiency) \\
Unified memory      & 16\,GB ($17.18\times10^{9}$ bytes); ${\sim}10$\,GB usable in practice \\
GPU backend         & Apple Metal (PyTorch MPS) \\
\midrule
\multicolumn{2}{l}{\textbf{Operating system}}\\
\midrule
macOS               & 27.0 (build 26A5388g), \code{macOS-27.0-arm64-arm-64bit} \\
\midrule
\multicolumn{2}{l}{\textbf{Software}}\\
\midrule
Python              & 3.12.12 (Homebrew; \emph{not} the system 3.14, which has no PyTorch wheels) \\
PyTorch             & 2.13.0, MPS built and available \\
NumPy               & 2.5.2 \\
gymnasium / Box2D   & 1.3.0 / 2.3.10 (CarRacing only) \\
cma                 & 4.4.4 \\
matplotlib / Pillow / PyYAML & 3.11.1 / 12.3.0 / 6.0.3 \\
\bottomrule
\end{tabular}
\caption{Execution environment. Recorded in every run's \code{bench.json} under
\code{meta.device}, so a stale result can always be traced to the machine that
produced it.}
\end{table}

\paragraph{Design decisions forced by the hardware.} The 10\,GB working budget
is not incidental; it determines the architecture of the pipeline.

\begin{itemize}[leftmargin=*]
\item \textbf{Frames are memory-mapped.} The frame store is the only large object
(${\sim}1$--$3$\,GB). Memory-mapping hands residency to the OS page cache, which
is right for the VAE's uniform random access; decompressing shards into the heap
is what pushes a laptop into swap.
\item \textbf{Latents are precomputed once.} This turns $M$-training from a
multi-gigabyte image problem into a ${\sim}30$\,MB vector problem, which is what
makes many seeds and many sequence lengths affordable at all.
\item \textbf{float32 throughout.} MPS mixed precision is unreliable for LSTM and
for the custom kernels in the Titans recurrence, and a numerics change landing on
only one arm would invalidate the comparison.
\item \textbf{CMA-ES workers run on CPU, not MPS.} Six Metal contexts cost more
in driver overhead and memory than they save for models this small, and it leaves
the GPU free.
\item \textbf{A resident-memory tripwire} aborts a run at 9.5\,GB with an
actionable message rather than letting the machine thrash.
\item \textbf{Data volumes were scaled down deliberately.} \cite{ha2018world} use
10{,}000 rollouts of 1000 steps---about 123\,GB of $64\times64$ frames. We use
$1200\times37$ for the primary configuration.
\end{itemize}

\paragraph{Measured cost.} Per-step and per-pass timings (batch 64, $L=37$):

\begin{table}[h]
\centering
\begin{tabular}{lrrr}
\toprule
Core & Backbone params & Training step & Inference step \\
\midrule
\code{lstm}    & $299{,}008$ & ${\sim}10$\,ms   & ${\sim}0.5$\,ms \\
\code{gru}     & $224{,}256$ & ${\sim}16$\,ms   & ${\sim}0.9$\,ms \\
\code{supreme} & $300{,}696$ & ${\sim}2000$\,ms & ${\sim}2.5$\,ms \\
\bottomrule
\end{tabular}
\caption{\textsc{Hope}'s two-order-of-magnitude training cost is inherent to a
per-timestep inner optimisation, not an artefact of a lazy implementation: the
chunk-parallel dual form of \cite[\S8.2]{behrouz2025nested} was implemented and
measured at $1.01\times$, $0.98\times$ and $1.05\times$ speedup at $L=37,69,250$
and then removed, because the cost is the backward pass through the sequential
recurrence rather than kernel launches. Both models see the same data for the
same number of epochs at matched parameter counts, so wall-clock must not be read
as ``spent more compute to learn more''.}
\end{table}

% =====================================================================
\section{Experimental protocol}

\begin{enumerate}[leftmargin=*]
\item Collect rollouts with a temporally correlated random policy. The dataset
directory is keyed by environment and collection settings, never by model, so
both arms read the same bytes.
\item Train $V$; freeze it; precompute $\mu,\log\sigma^2$ for every frame. The
sequence dataset resamples $z\sim\mathcal{N}(\mu,\sigma)$ each epoch, the
regularisation \cite{ha2018world} rely on.
\item Train $M$. Identical optimiser (Adam, cosine schedule), learning rate,
epochs, gradient clipping, batch size and seeds across arms; only the backbone
differs.
\item Train $C$ with CMA-ES, evaluating all candidates in a generation on the
same episode seeds (common random numbers).
\item Run the benchmark suite, writing one \code{bench.json} per (model, seed).
The reporting layer reads only those files, so any table can be regenerated
without retraining.
\end{enumerate}

\paragraph{A constraint that silently invalidates the study if violated.} A
training window must span cue$\rightarrow$decision. If \code{sequence\_length} is
shorter than that span, then no window contains both: windows starting at the cue
stop before the decision, and windows reaching the decision start after the cue.
The dependency is then absent from the training signal entirely and every memory
layer scores at chance for reasons unrelated to memory. Our first configurations
had exactly this defect. A guard now reads the collected data and aborts any run
that violates it.

% =====================================================================
\section{Results}
\label{sec:results}

All tables in this section are generated directly from the run artefacts by
\code{scripts/make\_tables.py} and pulled in with \code{\textbackslash input};
no number here is typed by hand, so the document cannot drift from the data.
Each cell is mean $\pm$ standard error over seeds, with the seed count printed
in the table so a single-seed entry is never mistaken for a converged one.
Live coverage is tracked in \code{docs/status.md}.

\IfFileExists{results_tables.tex}{% Auto-generated by scripts/make_tables.py -- do not edit by hand.
% Regenerate after new runs land; every number is a function of bench.json.

\begin{table}[h]
\centering
\small
\begin{tabular}{lrrrrr}
\toprule
Metric & \code{lstm} (MDN-RNN) & \code{gru} & \textbf{\code{supreme}} (Hope) & \code{supreme-cms} & \code{hope-attention} \\
\midrule
Decision accuracy $\uparrow$ & $0.475\!\pm\!0.047$ & $0.483\!\pm\!0.006$ & $1.000\!\pm\!0.000$ & $0.483\!\pm\!0.017$ & $1.000\!\pm\!0.000$ \\
NLL @ memory-critical $\downarrow$ & $25.55\!\pm\!2.15$ & $25.30\!\pm\!2.34$ & $15.61\!\pm\!0.97$ & $25.49\!\pm\!2.58$ & $16.73\!\pm\!0.33$ \\
NLL, all steps $\downarrow$ & $14.08\!\pm\!1.99$ & $14.12\!\pm\!1.97$ & $13.68\!\pm\!2.16$ & $13.90\!\pm\!2.06$ & $14.64\!\pm\!2.08$ \\
Effective memory horizon $\uparrow$ & $4.0\!\pm\!0.0$ & $20.0\!\pm\!4.0$ & $32.0\!\pm\!0.0$ & $4.0\!\pm\!0.0$ & $32.0\!\pm\!0.0$ \\
Probe AUC, state $\uparrow$ & $0.703\!\pm\!0.022$ & $0.890\!\pm\!0.018$ & $0.899\!\pm\!0.043$ & $0.699\!\pm\!0.010$ & $0.969\!\pm\!0.017$ \\
Return, real env $\uparrow$ & $0.062\!\pm\!0.187$ & $-0.125\!\pm\!0.125$ & 1.000 & -- & -- \\
Backbone params & 299,008 & 224,256 & 300,696 & 375,184 & 375,234 \\
\midrule
Seeds & 2 & 2 & 2 & 2 & 2 \\
\bottomrule
\end{tabular}
\caption{\code{TMaze-v0}, corridor 32 (a 32-step required memory horizon). Mean $\pm$ standard error over seeds. All columns share the same data, the same frozen $V$, the same MDN head, loss, optimiser and seeds; only the memory backbone differs.}
\label{tab:headline}
\end{table}

\begin{table}[h]
\centering
\small
\begin{tabular}{lcccccccc}
\toprule
Delay since cue & 0 & 1 & 2 & 4 & 8 & 16 & 24 & 32 \\
\midrule
\code{lstm} (MDN-RNN) & 1.00 & 1.00 & 1.00 & 0.96 & 0.43 & 0.40 & 0.42 & 0.41 \\
\code{gru} & 1.00 & 1.00 & 1.00 & 1.00 & 1.00 & 0.90 & 0.69 & 0.53 \\
\textbf{\code{supreme}} (Hope) & 1.00 & 1.00 & 1.00 & 0.53 & 0.97 & 0.82 & 0.89 & 0.98 \\
\code{supreme-cms} & 1.00 & 1.00 & 1.00 & 0.87 & 0.39 & 0.46 & 0.47 & 0.41 \\
\code{hope-attention} & 1.00 & 1.00 & 1.00 & 1.00 & 0.99 & 0.94 & 0.96 & 0.87 \\
\midrule
\emph{observation control} & 1.00 & 0.50 & 0.58 & 0.48 & 0.49 & 0.48 & 0.44 & 0.44 \\
\bottomrule
\end{tabular}
\caption{Linear probe accuracy of the cue from the recurrent state, by delay. The observation control probes $z_t$ alone: it is at ceiling while the cue is on screen and at chance from the next step onward, so any value above it is memory and nothing else.}
\label{tab:retention}
\end{table}

\begin{table}[h]
\centering
\small
\begin{tabular}{lrrrr}
\toprule
Metric & \code{lstm} (MDN-RNN) & \code{gru} & \code{attention} & \textbf{\code{supreme}} (Hope) \\
\midrule
Decision accuracy $\uparrow$ & -- & -- & -- & -- \\
NLL @ memory-critical $\downarrow$ & -- & -- & -- & -- \\
NLL, all steps $\downarrow$ & 10.87 & 11.13 & 10.19 & 10.36 \\
Effective memory horizon $\uparrow$ & 0.0 & 0.0 & 0.0 & 0.0 \\
Probe AUC, state $\uparrow$ & -- & -- & -- & -- \\
Return, real env $\uparrow$ & -- & -- & -- & -- \\
Backbone params & 300,032 & 225,024 & 299,314 & 300,952 \\
\midrule
Seeds & 1 & 1 & 1 & 1 \\
\bottomrule
\end{tabular}
\caption{\code{CarRacing-v0} \cite{ha2018world}. Reported for external comparability rather than as a discriminating benchmark: the task is close to fully observable from a single frame, so a memory layer has little to do.}
\label{tab:carracing}
\end{table}
}{%
\todo{results\_tables.tex not found --- run
\code{python scripts/make\_tables.py} in the repository.}}

\subsection{Average likelihood does not distinguish memory layers (final)}

On \code{TMaze-v0} corridor 24, an LSTM core and a GRU core---identical in every
other respect---reach one-step NLL of $12.87$ and $12.88$. Their effective memory
horizons are $4$ and $12$ steps. A study reporting only world-model likelihood
would have concluded that these two memory layers are the same.

\subsection{The probe control behaves as it must (final)}

\begin{table}[h]
\centering
\begin{tabular}{lcccccccc}
\toprule
Delay since cue & 0 & 1 & 2 & 4 & 8 & 12 & 16 & 24 \\
\midrule
GRU, state              & 1.00 & 1.00 & 1.00 & 1.00 & 1.00 & 0.80 & 0.48 & 0.48 \\
LSTM, state             & 1.00 & 1.00 & 1.00 & 0.68 & 0.48 & 0.48 & 0.48 & 0.48 \\
Observation control     & 1.00 & 0.48 & 0.48 & 0.48 & 0.44 & 0.48 & 0.48 & 0.48 \\
\bottomrule
\end{tabular}
\caption{Linear probe accuracy, \code{TMaze-v0} corridor 24. The control is at
$1.00$ while the cue is on screen and at chance from the very next step, which is
the sanity condition for the whole suite: anything above it is memory and nothing
else.}
\end{table}

\subsection{Retention and exploitation are separate milestones (final)}

At corridor 32 with a 60-epoch budget, \textsc{Hope} already \emph{held} the cue
(probe accuracy $0.867$ at delay 32, against $0.444$ for GRU and $0.311$ for
LSTM) but had not yet learned to \emph{use} it---all three models were at chance
on decision accuracy. The capability appears in the state well before it appears
in behaviour. A study stopping at the earlier budget would have reported ``no
difference in task performance'' while the mechanism was already in place. This
is a concrete methodological warning for anyone benchmarking memory layers, and
it is why both budgets are reported.

\subsection{Headline: a horizon neither recurrent core reaches}
\label{sec:headline}

Table~\ref{tab:headline} is the main result and Table~\ref{tab:retention} is the
mechanism behind it. At a 32-step required memory horizon, with matched parameter
counts and every other component shared, \textsc{Hope} reaches decision accuracy
$1.000\pm0.000$ while the MDN-RNN baseline ($0.475\pm0.047$) and the GRU control
($0.483\pm0.006$) both sit at chance. Likelihood on the memory-critical
transition separates by roughly ten nats---$15.61\pm0.97$ against $25.55\pm2.15$
and $25.30\pm2.34$---and the controller trained on \textsc{Hope}'s state solves
the maze outright (return $1.000$) where the controls do not.

The retention curve says why. The observation-only control is at ceiling while
the cue is on screen and at chance from the very next step, which is the sanity
condition for the whole suite: any value above it is memory and nothing else.
Against that baseline the cores separate cleanly---the LSTM's state has lost the
cue by delay~8 ($0.43$), the GRU's holds to delay~8 and decays to $0.53$ by
delay~32, and \textsc{Hope}'s is still at $0.98$ at delay~32. Effective memory
horizons of $4$, $20$ and $32$ steps.

Both \textsc{Hope} seeds agree \emph{exactly}---decision accuracy $1.000$ and
probe accuracy $0.978$ at delay~32 in both---which for a metric this close to a
step function is the reassuring outcome. And the advantage does \emph{not} appear
in average one-step NLL ($13.68\pm2.16$ against $14.08\pm1.99$ and
$14.12\pm1.97$, thoroughly overlapping): \S\ref{sec:findings-nll} restated with
final numbers, and the reason this study does not report likelihood alone.

\subsection{Ablation: the Continuum Memory System contributes nothing here}
\label{sec:ablation}

The ablation columns of Table~\ref{tab:headline} are more informative than the
headline, and they cut against the obvious reading of it.

\paragraph{CMS is inert.} \code{supreme-cms} is the \emph{baseline} LSTM with the
Continuum Memory System attached---25\% more parameters than the baseline---and
it lands exactly where the baseline landed: decision accuracy $0.483\pm0.017$
against $0.475\pm0.047$, memory horizon $4$ against $4$, memory-critical NLL
$25.49$ against $25.55$, probe AUC $0.699$ against $0.703$. Every one of those
differences is inside the seed noise. Whatever the multi-frequency knowledge
store does in the language-model regime it was designed for, at this scale and on
this task it does not contribute to retention at all.

\paragraph{The sequence mixer is what matters---and two different mixers both
work.} \code{hope-attention}, the paper's own variant with softmax attention in
place of the self-modifying Titans block, also reaches decision accuracy
$1.000\pm0.000$ and horizon $32$, with the best probe AUC of any core ($0.969$).
So the honest claim is narrower than ``self-modification is what remembers'':
\emph{the MDN-RNN's recurrent core is the bottleneck, and either a self-modifying
associative memory or attention removes it.} We cannot claim from these data that
self-modification is uniquely necessary.

\paragraph{But retention and imagination come apart, and that is the interesting
result.} \code{hope-attention} remembers essentially perfectly and yet imagines
badly: open-loop MSE at step~16 of $44.59\pm0.33$ against \textsc{Hope}'s
$0.506\pm0.17$, and a drift slope of $+0.126\pm0.002$ against \textsc{Hope}'s
$-0.070\pm0.013$---the only core in the study whose imagined trajectory diverges
rather than stays bounded. Attention caches every past frame and can always
retrieve the cue, but it does not learn dynamics that survive being fed their own
output. \textsc{Hope} does both, at a fixed-size state rather than one that grows
with the episode.

For a \emph{world} model that dissociation is the substantive finding. Retention
is necessary but not sufficient: the memory has to support rollout, and the
strongest retention baseline available is the worst open-loop predictor here.

\subsection{CarRacing: the external reference environment}

\code{CarRacing-v0} is the one environment we did not design, which is why it
matters for construct validity whatever it shows. What it shows is that it does
not discriminate: the task is close to fully observable from a single frame, so a
memory layer has little to do. One-step NLL across the four cores spans $10.19$
(attention) to $11.13$ (GRU), with \textsc{Hope} at $10.36$ and the MDN-RNN at
$10.87$---a range narrower than the seed-to-seed spread on \code{TMaze-v0}.
\textsc{Hope} has the best one-step reconstruction error ($0.873$ against
$1.085$--$1.118$) and the LSTM the best open-loop error at horizon~10; no
ordering is consistent across metrics.

We report this as a negative control, not as evidence for either model. All four
cores land within a nat of one another here while separating by ten nats on the
memory-critical transition of \code{TMaze-v0}: the discriminating power of the
synthetic suite is a property of the suite, not an artefact of the models.

\subsection{What these tables do not settle}

\begin{itemize}[leftmargin=*]
\item \textbf{Two seeds, not three.} With a near-bimodal metric, two seeds
agreeing exactly is encouraging but is not a distribution.
\item \textbf{One \textsc{Hope} run is short.} Seed~1's memory training was
interrupted at epoch~60 of~150. Its best-validation checkpoint was already past
the phase transition and reaches the same decision and probe accuracy as seed~0,
so it is reported---but the two seeds do not share a training budget and the
error bar on \textsc{Hope}'s NLL mixes two training lengths.
\item \textbf{\textsc{Hope}'s return is a single seed} (seed~0 only); no error bar.
\item \textbf{\code{supreme-titans} was not run.} Titans without CMS is
implemented and registered but never executed---at 55 minutes per seed it was cut
for time. Since CMS is demonstrably inert (above), we expect it to match
\textsc{Hope}, which would complete the $2\times2$ and upgrade ``CMS adds
nothing'' from an inference to a measurement. Until then that cell is empty.
\item \textbf{CarRacing ran at half scale} (150 rollouts, 12 memory epochs, no
controller).
\item \textbf{Three programme stages are outstanding:} the capacity axis
(\code{SequenceRecall-v0}), the task-incremental protocol, and the corridor
sweep---the last being what turns the choice of corridor~32 from a selected
operating point into a curve whose crossover is itself the finding.
\end{itemize}

% =====================================================================
\section{Threats to validity}

Stated plainly, because these are what a reader should press on.

\begin{enumerate}[leftmargin=*]
\item \textbf{The primary environment was built for this study.} A benchmark
designed alongside the model it evaluates is a construct-validity risk. The
independent checks---\code{SequenceRecall-v0} and especially
\code{CarRacing-v0}---are the mitigation.
\item \textbf{Operating-point selection.} Corridor 32 was chosen after observing
that corridor 16 saturates, with three of five cores at ceiling. Presented alone
this is post-hoc; presented as a sweep in which the crossover is the finding, it
is not. The corridor sweep is the mitigation and it is part of the programme.
\item \textbf{The headline requires $\lambda=20$.} Justified and measured
(\S\ref{sec:findings-nll}) and applied identically to both arms, but it is a
degree of freedom, and the $\lambda=1$ ablation must be reported at the operating
point rather than only at corridor 6.
\item \textbf{Decision accuracy is close to a step function.} A model either
learns the cue$\rightarrow$outcome rule or it does not, so per-seed values cluster
near $0.5$ or near $1.0$. The right summary is the mean \emph{and} the number of
seeds that learned the rule; a mean alone is misleading. Here both \textsc{Hope}
seeds learned it and none of the four control seeds did, which is the cleanest
form the statistic can take at $n=2$---but it is still $n=2$.

\item \textbf{The ablations are missing, and they carry the sharper claim.}
\textsc{Hope} is a composition of two mechanisms. Showing that the composition
beats an LSTM is a weaker result than showing which half is responsible, and the
latter is what \code{supreme-titans} and \code{supreme-cms} were built to
establish. They are implemented and unrun.
\item \textbf{Under-training is indistinguishable from forgetting} on the
decision metric, so budgets must be generous enough that the \emph{baseline} is
not merely undertrained---otherwise a win for the challenger is meaningless.
\item \textbf{Scale.} A $300$\,k-parameter memory on $64\times64$ synthetic
frames. Nothing here transfers automatically to the language-model regime the
nested-learning paper addresses.
\end{enumerate}

% =====================================================================
\section{Reproduction}

\begin{lstlisting}
cd "World Model/World Model"
./setup_env.sh && source .venv/bin/activate
make test                                    # 25 correctness tests
python wm.py study -c configs/tmaze_hard.yaml \
       --backbones lstm,gru,attention,supreme
python wm.py report runs
\end{lstlisting}

Each run writes \code{config.json}, \code{metrics.jsonl} and \code{bench.json};
tables and figures are regenerated from those artefacts alone.

% =====================================================================
\begin{thebibliography}{9}
\bibitem{ha2018world}
D.~Ha and J.~Schmidhuber.
\newblock World Models.
\newblock \emph{arXiv:1803.10122}, 2018.

\bibitem{behrouz2025nested}
A.~Behrouz, M.~Razaviyayn, P.~Zhong, and V.~Mirrokni.
\newblock Nested Learning: The Illusion of Deep Learning Architecture.
\newblock \emph{Advances in Neural Information Processing Systems (NeurIPS)},
2025. arXiv:2512.24695.

\bibitem{hansen2006cma}
N.~Hansen.
\newblock The CMA Evolution Strategy: A Comparing Review.
\newblock In \emph{Towards a New Evolutionary Computation}, 2006.

\bibitem{lopezpaz2017gem}
D.~Lopez-Paz and M.~Ranzato.
\newblock Gradient Episodic Memory for Continual Learning.
\newblock \emph{NeurIPS}, 2017.
\end{thebibliography}

\end{document}
